% !TeX program = xelatex
% ============================================================================
%  第 10 课　生成元、循环群与循环群的子群　—— 学生版讲义（原 12+13 合并）
%  与 02-课件/第10课-生成元与循环群/第10课-生成元与循环群.tex 同步
% ============================================================================

\xhlesson{10}{生成元、循环群与循环群的子群}{}

\begin{mubiao}
  \begin{itemize}
    \item 知道什么是生成元，能说出 \(D_4\) 里哪两个够用。
    \item 能找出模 \(12\) 加法里由 \(3\)、由 \(5\) 生成的子群。
    \item 会求模 \(12\) 各元素的阶，并发现子群大小都是 \(12\) 的因数。
  \end{itemize}
\end{mubiao}

\section*{一、生成元}

\begin{definition}[由 \(a\) 生成的子群]
  在群 \(G\) 里，把元素 \(a\) 和它反复做出来的所有结果放在一起，
  得到的那个小群，叫做\textbf{由 \(a\) 生成的子群}，记作 \(\langle a\rangle\)。
\end{definition}

\begin{definition}[生成元、循环群]
  如果整个群 \(G\) 能由一个元素 \(a\) 生成，也就是 \(G=\langle a\rangle\)，
  就说 \(a\) 是 \(G\) 的一个\textbf{生成元}，而 \(G\) 是一个\textbf{循环群}。
\end{definition}

\begin{dongshou}[走两遍圆盘]
  \begin{enumerate}
    \item 从 \(0\) 开始一直加 \(3\)：
      \(0\to3\to\underline{\hspace{1.4em}}\to\underline{\hspace{1.4em}}\to\underline{\hspace{1.4em}}\)（回到 \(0\)）。
      \quad \(\langle3\rangle=\)\underline{\hspace{5.4cm}}，一共 \underline{\hspace{1.2em}} 个。
    \item 从 \(0\) 开始一直加 \(5\)：先把它走的前几步补上——
      \[
        0\to5\to10\to\underline{\hspace{1.4em}}\to\underline{\hspace{1.4em}}\to\underline{\hspace{1.4em}}\to\cdots
      \]
      它一直走到重新回到 \(0\) 为止。数一数，一共走遍了 \underline{\hspace{1.2em}} 个数。
  \end{enumerate}
\end{dongshou}

\begin{sikao}
  为什么 \(3\) 只圈住 \(4\) 个，\(5\) 却圈住 \(12\) 个？\\
  （提示：看看它们和 \(12\) 的公因数。）
  \liukong{2}
\end{sikao}

\begin{example}
  \begin{center}
  \begin{tabular}{@{}lcc@{}}
    \toprule
     & 和 \(12\) 的公因数 & 圈住几个 \\
    \midrule
    \(3\) & \(1,3\)（最大的是 \(3\)） & \(12\div3=4\) \\
    \(5\) & 只有 \(1\) & \(12\div1=12\) \\
    \bottomrule
  \end{tabular}
  \end{center}
  公因数越大，圈住的越少；只有公因数 \(1\) 时，整整一圈都被圈住。
\end{example}

\section*{二、每个数的阶}

\begin{dongshou}[数几次回到 \(0\)，把阶填进表里]
  \begin{center}
  {\small
  \begin{tabular}{@{}l*{12}{c}@{}}
    \toprule
    \(a\) & 1 & 2 & 3 & 4 & 5 & 6 & 7 & 8 & 9 & 10 & 11 & 0 \\
    \midrule
    阶 & & & & & & & & & & & & \\
    \bottomrule
  \end{tabular}}
  \end{center}
\end{dongshou}

\begin{dingli}[阶与生成子群]
  元素 \(a\) 的阶，正好等于由它生成的子群 \(\langle a\rangle\) 的大小。
\end{dingli}

\begin{sikao}
  \(\langle4\rangle\) 和 \(\langle8\rangle\) 是同一个子群吗？说说理由。
  \liukong{2}
\end{sikao}

\section*{三、循环群里的所有小圈}

\begin{example}[模 \(12\) 加法里的全部子群]
  \begin{center}
  \begin{tabular}{@{}clc@{}}
    \toprule
    子群 & 是谁 & 大小 \\
    \midrule
    \(\{0\}\) & 只有 \(0\) & \(1\) \\
    \(\{0,6\}\) & 加 \(6\) 来回走 & \(2\) \\
    \(\{0,4,8\}\) & 由 \(4\) 生成 & \(3\) \\
    \(\{0,3,6,9\}\) & 由 \(3\) 生成 & \(4\) \\
    \(\{0,2,4,6,8,10\}\) & 由 \(2\) 生成 & \(6\) \\
    \(\{0,1,2,\dots,11\}\) & 全部十二个 & \(12\) \\
    \bottomrule
  \end{tabular}
  \end{center}
  一共就这 \(6\) 个。
\end{example}

\begin{sikao}
  最后一列是 \(1,2,3,4,6,12\)，和 \(12\) 有什么关系？\\[4pt]
  再想想：\(S_3\) 里是哪些数？\(D_4\) 里又是哪些数？三次的结论一样吗？
  \liukong{3}
\end{sikao}

\begin{xiaojie}
  用一句话说说：什么叫「生成元」，为什么要说「循环群」。
  \liukong{4}
\end{xiaojie}

\noindent\textbf{下节课}：回头去看幂的尾巴——\(2\) 的幂除以 \(7\)，余数会转圈。